\chapter{Sustainability and Social Commitment}
\label{ch:sustain}

This chapter records the sustainability dimension of the project,
following the matrix and reflection framework set out in the FIB
GEP sustainability guide~\cite{gepSustainabilityReport2018}.
\todo{FIB rubric expects explicit mapping to UN Sustainable Development Goals (e.g. SDG 4 Education, SDG 9 Industry/Innovation, SDG 11 Cultural Heritage, SDG 12 Responsible Consumption). Add a short paragraph or column in the matrix naming the SDGs each dimension addresses.}
\Cref{sec:sust-selfassessment} reports the self-assessment
that opened the project; \Cref{sec:sust-matrix} fills the
three-dimensional sustainability matrix (economic, environmental,
and social) at the project's design-and-development stage and at
its anticipated useful life; \Cref{sec:sust-commitment}
records the ethical position and open-science commitments adopted;
and \Cref{sec:sust-final} revisits the assessment in light
of the completed work.

% ─────────────────────────────────────────────────────────────────────────────
\section{Self-Assessment Summary}
\label{sec:sust-selfassessment}

After completing the GEP sustainability questionnaire, my strongest
area was identified as the \textbf{economic} dimension of software
projects, especially cost estimation, trade-off analysis, and
management control.
\todo{First-person voice throughout this section ("my strongest", "I am comfortable", "I generally", "I had not", "I also realised", "I now see"). Bachelor's thesis convention is impersonal/passive: rephrase to "the author's strongest area", "the questionnaire revealed", "this project's reflection process surfaced", etc.}
I am comfortable reasoning about whether a technical decision is
financially realistic within a constrained project, and about how
planning choices affect viability and long-term maintenance cost.
The questionnaire also showed a medium level in social reflection:
I generally consider accessibility and user benefit, but I still
tend to prioritise technical performance first and social impact
second when under time pressure.

The weakest area was the \textbf{environmental} lifecycle
perspective.
Before this project, I mainly associated environmental impact with
direct energy consumption during model training.
The reflection process made me consider additional factors that I
had not previously analysed in depth: the hardware manufacturing
footprint, end-of-life device management, and the difference
between one-time experimentation and sustained deployment impact.
I also realised that selecting lighter architectures is not only a
computational decision but also an environmental one---a point
explicitly reflected in the architectural justification of
\Cref{sec:mgmt-justification}.

Three concrete improvement commitments emerged for the remainder
of the project.

\textit{First}, energy-related impact should be quantified wherever
possible rather than asserted qualitatively.
The energy estimate of \Cref{sec:sust-environmental} is the
result of this commitment: it is computed from the actual GPU
hours consumed by the training runs, multiplied by the rated
power draw of the workstation and converted to CO\textsubscript{2}
using the official Spanish emission
factor~\cite{reeInformeSistemaElectrico2024,mitecoFactorEmision2025}.
This is a coarse calculation but it is grounded in measurable
quantities rather than in adjectives.

\textit{Second}, architecture choices should be justified with both
performance and efficiency criteria.
The comparison table in \Cref{sec:mgmt-justification} includes
``computational feasibility'' and ``data requirements'' as
first-class criteria alongside accuracy, and the backbone selected
in \Cref{sec:design-model} carries the environmental dividends
quantified in \Cref{sec:sust-environmental}.

\textit{Third}, the discussion of who benefits from the resulting
tool, and under what limitations, must be explicit rather than
implicit.
This commitment shapes the social-dimension discussion in
\Cref{sec:sust-social} and the ethical-position statement
in \Cref{sec:sust-commitment}: both name specific
beneficiary groups (musicians, educators, archivists,
researchers) and specific limitations (monophonic contract,
treble-clef-only inputs, residual SER profile) rather than
appealing to vague ``users'' and ``edge cases''.

The questionnaire was useful because it transformed sustainability
from a formal reporting requirement into a design constraint that
guides technical decisions throughout the project lifecycle.
It also strengthened my personal competency for future
professional work: I now see project quality as a balance between
technical success, economic feasibility, environmental
responsibility, and social value, and each of those four facets
needs to be reflected somewhere in the project plan rather than
treated as a downstream addendum.

% ─────────────────────────────────────────────────────────────────────────────
\section{Sustainability Matrix}
\label{sec:sust-matrix}

\Cref{tab:sustainability-matrix} summarises the
three-dimensional sustainability matrix as used in the GEP
framework.
The remainder of this section discusses each dimension in turn,
identifying the principal cost or benefit (PPP, ``the project as
designed and developed'') and the useful-life perspective (the
impact once the system is in service).

\begin{table}[h!]
    \centering
    \caption{Sustainability matrix for the project.
    PPP refers to the design-and-development stage; the useful-life
    column refers to deployment and downstream impact.}
    \todo{Caption is a full sentence with explanation. TFG style guide expects a noun-phrase caption with the long explanation moved to body text, e.g.\ ``Three-dimensional sustainability matrix (PPP vs.\ useful life).'' Same check for other captions in the chapter.}
    \label{tab:sustainability-matrix}
    \small
    \setlength{\tabcolsep}{5pt}
    \renewcommand{\arraystretch}{1.35}
    \begin{tabularx}{\linewidth}{@{} >{\raggedright\arraybackslash}p{2.3cm} L L @{}}
        \toprule
        \textbf{Dimension} & \textbf{PPP (design and development)} & \textbf{Useful life (deployment and impact)} \\
        \midrule
        Economic
            & Total project budget of \EUR{}16\,097.16, with explicit contingency and incidentals.
              The main cost driver is specialised engineering time.
            & Open-source release and fast local inference can reduce recurrent transcription cost compared with manual copying and commercial OMR tools. \\ \tfgrowsep
        Environmental
            & Estimated energy use of $\approx 57.5$\,kWh for the full project, mitigated by mixed-precision training, early stopping, and the reuse of existing
              hardware~\cite{reeInformeSistemaElectrico2024,mitecoFactorEmision2025}.
            & One-time training cost is amortised over repeated inference; local execution reduces cloud dependency for common use cases. \\ \tfgrowsep
        Social
            & The project strengthens the author's personal technical competencies and promotes responsible data and software practices.
            & Potential beneficiaries include musicians, educators, archivists, and researchers through improved accessibility, searchability, and preservation of music documents. \\
        \bottomrule
    \end{tabularx}
\end{table}

\subsection{Economic dimension}
\label{sec:sust-economic}

\textbf{PPP (design and development).}
The total project cost (\EUR{}16\,097.16; see
\Cref{sec:mgmt-budget}) is dominated by specialised
engineering time at roughly \SI{85}{\percent} of the total, a
profile typical of research-intensive software projects in which
labour outweighs material amortisation.
The full breakdown---salaries, hardware amortisation, indirects,
incidentals tied to the risk register, and the contingency
reserve---is constructed in \Cref{sec:mgmt-budget}.
\todo{Residual overlap with \Cref{sec:mgmt-budget}: the 85\% labour share is already stated in Ch.3. Either drop the percentage here and keep only the cross-reference, or recast this PPP cell as a sustainability \emph{interpretation} of the budget rather than a re-statement of it.}

\textbf{Useful life and state of the art.}
The current alternatives for lead-sheet digitisation are manual
transcription and general-purpose OMR tools, and the project's
economic improvement must be evaluated against both.
Manual transcription is expensive in recurrent cost (a trained
musician's time, at conservatory or freelance rates) and in
turnaround time (typically of the order of minutes to tens of
minutes per page, depending on density and clarity).
General-purpose OMR tools---both commercial (SmartScore,
PhotoScore, SharpEye) and open-source (Audiveris)---reduce the
recurrent cost but require substantial correction effort on the
jazz-specific notation studied here, because their internal
heuristics are tuned for classical engraving styles
(\Cref{sec:soa-classical}).
A specialised CRNN--CTC model targeted at the Real Book
engraving style reduces the correction workload further, on the
class of input it was trained for, and eliminates the licence
costs that the commercial alternatives charge per seat.

\textbf{Improvement contribution.}
The principal economic improvement is not zero-cost digitisation
but a \textit{lower cost per corrected page} after deployment.
A user of the system still has to verify and correct the output
manually for any application that depends on note-perfect
transcription, but the per-page correction time is reduced
because the model handles the bulk of the symbol recognition
correctly (\Cref{ch:results}).
A second improvement is research reproducibility: the open-source
release of the code, the training scripts, and the trained model
weights lowers the entry cost for future OMR research targeting
the same or adjacent domains, because the synthetic data pipeline
and the CRNN--CTC training stack do not have to be rebuilt from
scratch.

\subsection{Environmental dimension}
\label{sec:sust-environmental}

\textbf{PPP (design and development).}
The project's estimated energy consumption is approximately
\SI{57.5}{\kilo\watt\hour}, derived from the rated power draw of
the development workstation (a desktop with an NVIDIA RTX~3060
GPU and a typical full-system draw of \SI{250}{\watt} during
training) multiplied by the estimated cumulative training and
data-generation time over the project's twenty-one-week schedule.
\todo{Methodology lacks the actual numbers behind 57.5\,kWh. State the cumulative GPU-hours (training + data generation + inference) and show the multiplication explicitly, e.g. $230\,\text{h} \times 0.25\,\text{kW} = 57.5\,\text{kWh}$, so the figure is reproducible from the cited training logs.}
At the Spanish electricity emission factor of approximately
\SI{0.23}{\kilo\gram\per\kilo\watt\hour} of CO\textsubscript{2}
reported by Red El\'ectrica de Espa\~na and
MITECO~\cite{reeInformeSistemaElectrico2024,mitecoFactorEmision2025},
this corresponds to roughly \SI{13.2}{\kilo\gram} of
CO\textsubscript{2} emitted during the project's execution.
The estimate is coarse but is grounded in measurable quantities
rather than rules of thumb.

The dominant contributor to the total is GPU training time.
Inference, dataset generation, and development sessions away from
the GPU contribute much less per hour because they do not engage
the high-power-draw components of the workstation.
This concentration of environmental impact in a single activity
makes the impact tractable: optimising training has substantially
more environmental leverage than optimising the rest of the
workflow.

\textbf{Minimisation plan.}
Four mitigations together keep the training-time energy cost
within the estimated envelope.
First, the architecture choice (\Cref{sec:design-model})
favours a \textit{lightweight} backbone, ResNet18, over the
heavier alternatives that would yield only marginal accuracy
gains on monophonic input;
\Cref{sec:mgmt-justification} makes this trade-off
explicit.
Second, \textit{mixed-precision training} via PyTorch's
\code{torch.amp} reduces both VRAM occupancy and the cycle time
per training step, which translates directly into fewer
GPU-hours per converged model.
Third, \textit{early stopping} with a twelve-epoch patience halts
training once the validation SER stops improving, avoiding the
common pattern of running for a fixed number of epochs ``to be
safe''.
Fourth, the project \textit{reuses existing hardware} rather than
provisioning new equipment specifically for the thesis; the
embodied energy of manufacturing a new workstation would dwarf
the operational energy of the project by an order of magnitude.

Data generation and training are also organised in batches to
avoid redundant runs: the dataset is fixed before model training
begins, so that data-pipeline iteration and model iteration do
not compound in cost.
Each significant training run is logged and its CSV summary
retained, which makes it possible to identify and avoid
duplicate experiments in later sprints.

\textbf{Useful life and improvement contribution.}
Once trained, inference has a low marginal energy cost compared
with retraining: a single page transcription completes in
seconds on consumer hardware and consumes a small fraction of a
watt-hour of energy.
The energy break-even point of the project---the number of
inference invocations at which the operational energy of inference
equals the one-time training cost---is therefore high enough that
the trained model can be regarded as an amortised asset.
Sharing the trained weights and the training scripts spreads the
one-time training cost across multiple downstream users and
multiple academic iterations, so that the per-application share
of the project's environmental footprint shrinks as the model is
re-used.

\subsection{Social dimension}
\label{sec:sust-social}

\textbf{PPP (design and development).}
In personal-development terms the project strengthens
competencies in three areas relevant to the author's future
professional work: machine-learning engineering (architecture
design, training-loop implementation, evaluation methodology),
project planning (work-breakdown structures, sprint cadence,
budget management), and scientific communication (technical
writing, figure-driven argument, and the citation discipline
required by a thesis-length document).
The project also reinforces ethical awareness about the
limitations and failure modes of automated transcription:
working through the residual-error analysis of
\Cref{ch:results} produces a much more concrete picture
of what the model gets right and wrong than a single aggregate
metric can convey.

\textbf{Useful life and real need.}
There is a real demand for digitising lead sheets used in music
education and performance, and the system is built specifically
to serve that demand on a notation style (Real Book engraving)
that is widely used and poorly served by existing tools.

For \textit{musicians}, the downstream applications include
automatic transposition to a different key (essential for
B$\flat$, E$\flat$, and F transposing instruments),
MIDI playback for practice and accompaniment, and integration
with practice apps that depend on a symbolic representation of
the score.
For \textit{teachers}, the value is in producing digital
materials that can be edited, transposed, and annotated for
specific pedagogical purposes without re-typesetting the entire
piece by hand.
For \textit{archivists}, OMR is a precondition for making large
paper collections searchable: a scanned PDF is preserved but not
discoverable, whereas a symbolic representation can be queried
melodically, harmonically, or rhythmically.
For \textit{researchers}, especially in music information
retrieval and computational musicology, a symbolic version of a
large jazz corpus enables statistical and structural analyses
that simply cannot be performed on raster images.

Across all four groups, the project can improve quality of life
in two related ways: by reducing the repetitive manual work
required to convert paper notation to a machine-readable form,
and by facilitating the long-term preservation of music content
that currently exists only as paper or low-resolution scanned
PDF.
\todo{Social dimension only names beneficiaries; it does not name groups \emph{excluded} or potentially harmed. Add at least one sentence on users the current system does not serve (e.g.\ players of non-Western notation, full-score/polyphonic users, visually impaired musicians needing audio-first output, users of handwritten or pre-engraved historical sources). The FIB rubric expects both sides.}

\textbf{Boundaries.}
Social value depends on transparent communication of the
system's current limitations.
Four limitations are particularly important to name explicitly.
The \textit{monophonic-only contract}
(\Cref{sec:design-data}) means the model is not designed
to transcribe piano scores, choral parts, or any input with more
than one simultaneous voice on the staff.
The \textit{Real-Book-specific engraving style} means the model's
accuracy on other engraving styles---classical Sibelius output,
historical editions, handwritten manuscripts---is not the
accuracy reported in \Cref{ch:results} and would have to
be measured separately.
The \textit{treble-clef-only assumption} excludes any input on
F or C clefs, which is acceptable for lead sheets but
disqualifying for full scores.
The \textit{residual error rate} reported in
\Cref{ch:results}, while small in aggregate, is not zero;
applications that depend on note-perfect transcription must
include a human verification step.
Overpromising the system in scenarios outside the validated
scope would erode rather than enhance its social value, and
would also expose the system's users to downstream costs they
have not been warned about.

% ─────────────────────────────────────────────────────────────────────────────
\section{Commitment and Ethical Position}
\label{sec:sust-commitment}

The project's ethical position aligns with the open-science
commitments of the GEP sustainability framework, and is recorded
here as five concrete practices rather than a single statement of
principle.

\textit{Open source release.}
The source code is released under a permissive open-source
licence and is tracked in a public Git repository.
Every experiment in the thesis can be reproduced from a clean
clone using a pinned-dependency environment managed by Poetry
(\Cref{sec:impl-reproducibility}), and the trained model
weights are distributed alongside the code.
This commitment matters because OMR systems often disappear into
private repositories that cannot be inspected or extended; making
the implementation public is the most direct contribution the
project can make to the surrounding research community.

\textit{Documented evaluation protocol.}
The metric protocol used in \Cref{ch:results} is documented
in full, including the definition of the aggregate Symbol Error
Rate, the melodic SER that strips structural tokens, and the
category-level error breakdown that decomposes edits into
substitutions, insertions, and deletions per token class.
Future work can therefore replicate the present evaluation
exactly, or extend it with additional metrics, without having to
guess at how the present numbers were computed.

\textit{Licence-respecting use of third-party data.}
All third-party data used in the project is consumed in
accordance with its licence.
PrIMuS, the synthetic-supervision source, is publicly distributed
for research use~\cite{PrIMuSDataset} and is used here unmodified.
The author's own scans of the Real Book are used only for
fine-tuning and evaluation in a strictly academic context; the
trained weights cannot be redistributed in a form that recovers
the underlying copyrighted material.
\todo{Ethics/IP section is too thin for the FIB rubric. (1) Name the exact PrIMuS licence (CC-BY-NC-SA? research-use-only?) rather than ``publicly distributed for research use''. (2) State explicitly that the Real Book is a copyrighted commercial publication (Hal Leonard) and explain why training on private scans is defensible under academic fair use / TDM exceptions (e.g.\ Spanish LPI art.~67 / EU DSM Directive art.~3). (3) Confirm the released weights and evaluation set do not embed reconstructable copyrighted pages.}

\textit{No personal data, no safety-critical decisions.}
The system does not process personal data at any stage---input
images are music notation, output sequences are symbolic
representations of that notation---and it is not deployed for
safety-critical automated decision-making.
This narrows the ethical surface of the project to its
intellectual-property dimension and to its honesty about its
capabilities, which is the subject of the next practice.
\todo{Add a short \emph{misuse cases} paragraph: e.g.\ users could digitise and redistribute copyrighted lead sheets at scale, bypass paid score libraries, or generate derivative works without rights-holder consent. Note how the open-source release mitigates (transparency, opt-in) versus exacerbates (lower barrier) each risk. The FIB ethics rubric expects misuse to be named, not only personal-data and safety.}

\textit{Documented limitations.}
The known limitations of the system---the monophonic contract,
the treble-clef-only assumption, the eight allowed time
signatures, and the residual SER profile---are documented
prominently in the project's \code{README}, in the project
overview documentation, and in the present thesis
(\Cref{sec:scope} and \Cref{sec:sust-social}).
Future users of the model can therefore evaluate in advance
whether the system is suitable for their use case rather than
discovering its limitations through failure in production.

% ─────────────────────────────────────────────────────────────────────────────
\section{Final Reflection}
\label{sec:sust-final}

Revisiting the sustainability assessment at the end of the project,
two updates are worth recording.

First, the \textbf{environmental} estimate of $\approx 57.5$\,kWh
made at the planning stage is consistent with the executed
training profile: the final checkpoint converged in
\num{37} epochs (\Cref{sec:res-training}) and the full
training run including ablations fits within the planned envelope.
Mixed-precision training was the single most effective mitigation:
it reduced VRAM usage enough that the entire pipeline runs on a
consumer GPU, avoiding the cloud-rental scenario reserved as the
worst-case fallback in the incidentals table
(\Cref{tab:incidentals}).

Second, the \textbf{social} value of the system depends on the
deployment surface, not just on its accuracy in isolation.
The FastAPI web service of \Cref{sec:impl-api} packages
the inference pipeline behind a single REST endpoint with no
GPU requirement on the client side, which is precisely the
deployment shape that makes the tool accessible to musicians,
educators, and archivists who do not have local ML expertise.
\todo{Unhedged future claim: ``makes the tool accessible to musicians, educators, and archivists''. The API is local-only and undeployed in this thesis; soften to ``is designed to be accessible'' or ``may make the tool accessible, subject to a hosted deployment''. Also check the chapter for other unhedged claims (``the project can improve quality of life'', ``Future users\dots\ can therefore evaluate'') and switch to ``may'' / ``is intended to''.}
This shape was a planned design decision, but its sustainability
significance only becomes visible once the technical work is
complete: a model with the same accuracy and no accessible
deployment surface would have substantially less social impact
regardless of its training-time efficiency.

A residual concern, recorded here for completeness rather than as
an unresolved issue, is the \textbf{lifecycle} of the trained model
weights themselves.
\todo{The Environmental dimension covers operational kWh and embodied hardware energy, but says nothing about end-of-life: GPU disposal/e-waste, model-weight archival storage cost, and the energy cost of \emph{retraining} when the model becomes obsolete. Add one or two sentences here or in \Cref{sec:sust-environmental} so the lifecycle perspective is complete (manufacture $\to$ operation $\to$ end-of-life).}
Model files have a finite useful life---they become obsolete as
new architectures are released, as datasets evolve, and as the
target notation style drifts.
Publishing the training scripts alongside the weights is the
standard mitigation: it ensures that the cost of training a
replacement model can be paid once and then redistributed, rather
than being repeated independently by each successor project.
